\SingleSpacing

\chapter{Abstract}

\OnehalfSpacing

Approximately 25 \% to 40 \% of anthropogenic carbon dioxide ($\mathrm{CO_2}$) emissions have been absorbed into the oceans, influencing the physics and chemistry of seawater and contributing to ocean warming (OW) and acidification (OA).
% These environmental stressors interact to affect biochemical reaction rates and physiological functioning in marine biota, thus expediting or mitigating biological impacts.
Marine organisms will have to adjust their physiology and behaviour to cope with these changes and increase their resilience when exposed to these stressors.
% However, whether they are able to maintain their fitness in the face of future climate change is dependent on the severity and time scale of the stressors involved.
This thesis investigates the impact of long-term exposure to OW and OA, at levels predicted for up to the year 2100 on the physiological responses of the sea urchin, \textit{Heliocidaris erythrogramma}. 
This species is an ecologically important benthic ecosystem engineer in south-east Australia with the ability to reshape community assemblages through its herbivory.
Experimental studies on adult \textit{H. erythrogramma} were conducted using experimental designs which involved gradual introduction to future OW or OA conditions over long-term incubation periods to accomodate acclimatory responses, and temperature adjustments to incorporate seasonal change. The influence of parental acclimation effects on offspring responses to OW and OA was also investigated. 

Temperature is the most important factor determining the biological performance and fitness of marine ectotherms, and in the southeast Australian OW hotspot, up to +6 \degree C of warming is projected for the region by the year 2100. 
The consequences of exposure to warming was investigated in \textit{H. erythrogramma} through gradual adjustment over 6 weeks to elevated temperature (ambient 20 \degree C, 22 \degree C, 24 \degree C, 26 \degree C) followed by an acclimation period of 3 months at these temperatures in constant conditions.
Metabolic rate, $Q_{10}$, heat-shock protein (HSP70) expression, gonad index, gonad histology and survival were measured. 
There were no significant effects of elevated temperature on metabolic rate, HSP70 protein expression and survival at 22 \degree C, indicating that the sea urchin was able to tolerate a +2 \degree C increase in temperature.
However, metabolic depression was evident at the highest warming level of 26 \degree C which caused decreased survival and elevated HSP70 protein expression.
\textit{Heliocidaris erythrogramma} responded to the stress of increased temperature by altering its physiology, but long-term exposure to +4 \degree C and +6 \degree C levels of warming incurred unsustainable metabolic costs and led to trade-offs between metabolism, HSP70 expression and survival. 

Decreased pH through uptake of $\mathrm{CO_2}$ affects a wide range of physiological functions in sea urchins due to alteration of the carbonate chemistry of seawater which affects pH balance, mineral saturation for skeletogenesis and ongoing hypercapnia (\textit{p}$\mathrm{CO_2}$).
The thermal and metabolic responses of \textit{H. erythrogramma} are likely to be influenced by OA, with the combined impact of both OW and OA likely to be difficult to predict.
To investigate the effects of both stressors on the physiology of \textit{H. erythrogramma}, OA was incorporated with OW in a multistressor experimental design.
\textit{Heliocidaris erythrogramma} collected in winter were gradually adjusted over 7 weeks to temperature (ambient, +2 \degree C and +3 \degree C) and pH (ambient 8.0, low 7.6) treatments.
Temperature was adjusted weekly to reflect the seasonal temperature change that occurs over time from winter to summer (17 \degree C to 22 \degree C). 
Thus the elevated temperature treatments were adjusted to offset the temperature profile.
The sea urchins were incubated in those conditions for 22 weeks.

Metabolic rate was quantified at 4 and 12 weeks of acclimation and feeding and ammonia excretion rates were measured at 12 weeks. 
Survivorship was recorded daily for up to 22 weeks.
In single stressor treatments metabolic rate increased with either OW or OA.
In contrast, the two stressors interacted when combined causing a decrease in metabolism as temperature increases at low pH.
There was no evidence of interaction between the two stressors for feeding rate, ammonia excretion rate, or survival. 
Thus, while pH influenced metabolic rate causing metabolic depression at +3 \degree C, there was no evidence of pH affecting the other parameters measured. 
Survival was significantly reduced at +3 \degree C, regardless of pH, but not at +2 \degree C. 
% It appeared that the earlier arrival of warmer temperature at +3 \degree C was stressful, and the effect of this 1 \degree C difference was evident due to the inclusion of seasonal variation in the experiment.
It appeared that the earlier arrival of warmer temperature in the +3 \degree C treatment was stressful, especially during a time of increasing gonad development. 
It appears that \textit{H. erythrogramma} may be living close to its lethal limit in Sydney Harbour since it already experiences the temperatures used in the study. Current populations are likely to be vunerable to the combined stress of OW and OA in the near future.

While investigations of climate change effects at the adult stage provides insights into the resilience of established populations, it is also important to understand how their offspring may be affected.
% Long-term studies may provide crucial information on the impacts of climate change stressors on the resilience or current populations of organisms, but the potential for adaptation (or acclimation) over time may also require an understanding of how responses may be reshaped over generations.
These transgenerational carry-over effects are influenced by a combination of parental environmental history and developmental plasticity, and can convey resilience to offspring between generations.
To investigate the influence of parental environment on offspring performance, adult \textit{H. erythrogramma} were acclimated to OW and OA conditions over the gonad developmental period (winter to summer), and then spawned after 3 months in treatments.
Embryos were incubated over 14 days to the juvenile stage under a range of temperatures (18 \degree C, 20 \degree C, 22 \degree C, 24 \degree C) and two $\mathrm{pH_T}$ levels (8.0, 7.6).
Juvenile metabolic rate and test diameter (as a proxy for growth) were measured.
It was found that parental exposure to warming had a positive effect on the metabolic rate responses of juveniles to temperature, whereby they respired more (up to 30 \%) when incubated in warmer treatments.
There were no effects of parental environment or juvenile thermal environment on growth.
Parental exposure to low pH had no effect on metabolic rate or growth, however exposure to low pH during offspring development significantly affected growth of the juveniles (up to 4 \% decrease).

% All of the studies performed in this thesis used metabolic rate as a physiological indicator of performance. 
This thesis uses metabolic rate as a key indicator of physiologal response to projected climate change conditions.
Metabolic rate measurements are well-established as a key parameter to assess the biological and ecological consequences of environmental stress. 
To address problems associated with the lack of standard methods to analyse increasingly complex respirometry data, a new package for the R statistical computing language, called \texttt{respR}, was developed.
This software package provides a range of functions for analying respirometry data.
Importantly, this package simplifies handling and analysis of respirometry data and can be used to reliably and rapidly generate reproducible analyses of metabolic rates.
The \texttt{respR} package provides new methods to accurately detect maximum and minimum rates, linear sections of previously intractable data, and critical oxygen tension, based novel techniques of rolling regression and machine learning.
The package has been published as an open-source project to address data reproducibility and transparency when processing respirometry data.


% This thesis provides evidence that the physiological and biological responses of animals to climate change stressors are influenced by the interactions between stressors, seasonal changes and trans-generational effects.
For \textit{H. erythrogramma}, an ecologically important species on Australia's temperate reefs, there were clear long-term physiological consequences of acclimation to OW and OA on physiological condition and fitness. 
This species is vunerable to a +3 \degree C increase in temperature, but appears resilient to decreased pH.
The response of \textit{H. erythrogramma} to the combined effect of both stressors appeared to be dominated by changes in temperature.
The outcomes of this research have significant implications as to the resilience of \textit{H. erythrogramma} to climate warming, especially where it resides in the global warming hotspot of temperate Australia.
However, there was evidence of carry-over effects on juvenile physiology shaped by parental environment, which suggests a capacity for transgenerational acclimation and resilience. 
Further work is required to determine the multi-generational response of this species to climate change.
% Lastly, data reproducibility was addressed with the introduction of an R package for respirometry which will allow physiologists to process respirometry data transparently.


% Overall this thesis shows that the sea urchin \textit{H. erythrogramma} was vunerable to near future OW and OA, but evidence of carry-over effects on juvenile metabolism, shaped by parental environment, suggests that further work is required to determine the resilience of this species to climate change. 
% In addition, a novel computing package was developed which will hopefully allow future users to handle and process extensive respirometry datasets in a manner that is fully transparent and reproducible.


\cleardoublepage

